%!TEX root = TheoretischePhysikC.tex
% Skript v0.3, Kapitel 1: Kern von C0 und C1; Quellen in docs/plan/script-v0.3/sync-ch1.md.

\section{Vorlesung: Die Geburt der Quantenmechanik}
Wir befinden uns in einem Quanten-Universum. Bis 1900 haben wir die Quantenaspekte des Universums nicht bemerkt, weil uns die klassische Physik eine sehr gute Annäherung an die Physik des Alltags geliefert hat. Erklärungsbedürftig ist also nicht die Quantenmechanik, sondern die Güte der klassischen Näherung. Das Stichwort dazu heißt \emph{Dekohärenz}; sie gehört in die fortgeschrittene Quantenmechanik.

Eine Bemerkung zur Methode, die für den ganzen Kurs gilt: Man kann erst die falschen, klassischen oder halbklassischen Modelle durchrechnen, um dann zu zeigen, dass sie nicht funktionieren. Das ist nicht mein Stil. Wir betrachten von Anfang an alles richtig; das Bohrsche Atommodell rechnen wir nirgends durch. Dieses Skript enthält nur den Kern; Herleitungen, Vertiefungen und interaktive Abbildungen stehen im Webkurs.

\subsection{Warum eine neue Theorie nötig war}
Die Vorgeschichte der Quantenmechanik ist eine Reihe von Beobachtungen, die die These, unser Universum werde durch die klassische Physik beschrieben, allmählich in Frage stellten.
\begin{enumerate}
	\item \emph{Schwarzkörperstrahlung} (Planck, 1900). Ein Schwarzkörper reflektiert keine einfallende Strahlung, strahlt bei einer Temperatur $T$ aber dennoch. Das Spektrum dieser Strahlung ist mit klassischer Physik nicht erklärbar. Das richtige liefert die Annahme, dass Licht der Frequenz $\nu$ aus \emph{Lichtquanten} der Energie $E = h\nu$ besteht. Planck hat die Konstante $h$ ursprünglich als ,,Hilfsgröße`` postuliert, um das Strahlungsgesetz herzuleiten:
	\begin{equation*}
		u(\nu, T)\,\dd\nu = \frac{8\pi h \nu^3}{c^3}\,\frac{\dd\nu}{e^{h\nu/(\kB T)} - 1}
	\end{equation*}
	ist die Energie pro Volumen im Frequenzintervall $\dd\nu$, mit der Lichtgeschwindigkeit $c$ und der Boltzmann-Konstante $\kB$. Dieses Gesetz leiten wir im statistischen Teil selbst her. In der theoretischen Physik verwendet man vorwiegend die \emph{reduzierte Planck-Konstante}
	\begin{equation*}
		\eqbox{\hbar = \frac{h}{2\pi}}
	\end{equation*}
	\item \emph{Photoeffekt} (Lenard; Einstein ab 1905). Licht kann Elektronen auslösen. Bei festem Material hängt ihre \emph{maximale} kinetische Energie von der Lichtfrequenz ab, nicht von der Intensität. Klassisch ist das merkwürdig: Sonst bedeutet mehr Intensität mehr Energie. Gibt ein Lichtquant seine Energie an ein einzelnes Elektron ab, so lautet die Bilanz
	\begin{equation*}
		E_{\mathrm{kin,max}} = \tfrac12 m_{\mathrm{e}} v_{\max}^2 = h\nu - W
	\end{equation*}
	mit der Bindungsenergie $W$ des Materials (nichtrelativistisch geschrieben). Unterhalb der Frequenz $W/h$ löst in diesem Modell ein einzelnes Lichtquant kein Elektron aus.
	\item \emph{Stabilität der Atome, Linienspektren, Franck-Hertz-Versuch}. Nach den Liénard-Wiechert-Potentialen strahlt jede beschleunigte Punktladung. Ein Elektron auf einer Bahn um den Kern müsste also ständig Energie verlieren und in den Kern stürzen, was es nicht tut. Linienspektren und der Franck-Hertz-Versuch (1914) deuten darauf hin, dass in Atomen nur bestimmte, diskrete Energieniveaus erlaubt sind.
	\item \emph{Compton-Effekt}. Wie sich bei der Streuung von Licht an Elektronen Energie und Impuls verteilen, beschreibt man richtig mit Lichtquanten der Energie $h\nu$ und des Impulses $p = h/\lambda$ und mit Energie- und Impulserhaltung.
	\item \emph{Stern-Gerlach-Versuch} (1922). Ein Strahl von Silberatomen läuft durch ein inhomogenes Magnetfeld. Klassisch sollte er sich kontinuierlich auffächern; beobachtet werden zwei getrennte Strahlen. Bei fester Drehimpulsquantenzahl nimmt die Komponente des Drehimpulses in Feldrichtung nur endlich viele Werte an. Aus solchen quantisierten Drehimpulsen kann man heute Qubits bauen.
	\item \emph{Interferenz bei Elektronen}. Die Beugung von Elektronen an Kristallen (Davisson und Germer, 1927) war ein Hinweis auf Welleneigenschaften des Elektrons, die de Broglie 1924 postuliert hatte.
\end{enumerate}
Ich fasse die Befunde in zwei Sätzen zusammen.

\begin{merksatz}
Materie und Strahlung besitzen sowohl Wellen- als auch Teilcheneigenschaften.
\end{merksatz}

\begin{merksatz}[Die wichtigste Aussage über die Quantenmechanik]
Die Quantenmechanik ist eine \emph{probabilistische Theorie}. Sie liefert uns somit nur Vorhersagen in Form von \emph{Wahrscheinlichkeiten}.
\end{merksatz}

Diese zweite Aussage ist der ersten übergeordnet: ,,Wellen- und Teilcheneigenschaften`` bleibt ein wenig schwammig, die probabilistische Aussage dagegen ist präzise und liegt allen Quantenphänomenen zugrunde. Klassisch lauten die Vorhersagen definitiv: Das Teilchen ist zu dieser Zeit an dieser Stelle. Wahrscheinlichkeiten treten dort allein wegen fehlender Information auf; fundamentale Wahrscheinlichkeiten gibt es in der klassischen Physik nicht. Die statistische Deutung hat Born 1926 geliefert, und sie ist für mich der wichtigste Aspekt der Theorie. Nehmen Sie den Satz bitte mit: In jeder mündlichen Prüfung frage ich nach der wichtigsten Aussage über die Quantenmechanik.

Die Interpretation der Quantenmechanik ist trotz vieler Ideen bis heute umstritten; einen Widerspruch zu ihr gibt es bis heute nicht. Unsere Annahme ist, dass die Quantenmechanik die korrekte Beschreibung des Universums ist.

\subsection{Was eine physikalische Theorie liefert}
Aus dem zweiten Merksatz mache ich einen Grundsatz für jede physikalische Theorie: Physikalische Theorien sind probabilistische Theorien, ihre Vorhersagen haben die Form von Wahrscheinlichkeiten. Deterministische Theorien fasse ich als Teilmenge der probabilistischen auf. Eine physikalische Theorie liefert Daten für Präparationen, Dynamik und Messung.

\begin{greybox}[Jedes Experiment in drei Schritten]
Für den Theoretiker sehen alle Experimente der Physik gleich aus: Etwas wird \emph{präpariert}, etwa ein Würfel oder im Labor ein Lichtstrahl; das Experiment wird \emph{durchgeführt}; am Ende \emph{messen} wir und erhalten eines der möglichen Ergebnisse.
% fig_vl01_experiment.tex — Jedes Experiment in drei Stadien: Präparation,
% Durchführung, Messung; die Klammer fasst die ersten beiden zu einer
% komplizierteren Präparation zusammen.
%
% Adaptiert aus content/c00-theorie/figures/experiment.tex (c00.fig.experiment),
% das wiederum den grauen Kasten ,,Jedes Experiment in drei Schritten``
% des Skripts v0.2 (vl01.tex:118-126) zeichnet.  Inhalt unverändert; nur die
% Stile sind die des Skripts (tpc-style.sty, tpc-macros.sty) statt der
% Design-Token des Kurses (tpc state, tpc dynamics, tpc effect), damit die
% Abbildung mit scripts/check_lecture.sh allein baut.
%
% Geometrie: ein Schema, keine Kurve.  Drei gleich große Kästen der Breite
% w = 2.9 und Höhe 1.3, getrennt durch Pfeile der Länge g = 0.9; der i-te
% Kasten (i = 0, 1, 2) beginnt bei x_i = i (w + g), also bei 0, 3.8, 7.6.
% Die Klammer reicht von x_0 = 0 bis x_1 + w = 6.7.
% Rollen wie im Kurs: Präparation durchgezogen, Durchführung mit dickerem
% Strich, Messung gestrichelt (graustufentauglich, Farbe nur zusätzlich).

\begin{center}
\begin{tikzpicture}[x=1cm, y=1cm]
  \def\w{2.9}\def\g{0.9}\def\h{1.3}
  % 1 Präparation: durchgezogen
  \draw[cgblue, line width=0.9pt] (0,0) rectangle (\w,\h);
  \node[lbl, align=center] at (0.5*\w,0.5*\h)
    {1\\[1pt]\tpcifenglish{preparation}{Präparation}};
  \draw[axisstyle] (\w,0.5*\h) -- (\w+\g,0.5*\h);
  % 2 Durchführung: dickerer Strich
  \draw[spacecadet, line width=1.6pt] (\w+\g,0) rectangle (2*\w+\g,\h);
  \node[lbl, align=center] at (1.5*\w+\g,0.5*\h)
    {2\\[1pt]\tpcifenglish{run}{Durchführung}};
  \draw[axisstyle] (2*\w+\g,0.5*\h) -- (2*\w+2*\g,0.5*\h);
  % 3 Messung: gestrichelt
  \draw[birrenorange, line width=0.9pt, dashed]
    (2*\w+2*\g,0) rectangle (3*\w+2*\g,\h);
  \node[lbl, align=center] at (2.5*\w+2*\g,0.5*\h)
    {3\\[1pt]\tpcifenglish{measurement}{Messung}};
  \draw[axisstyle] (3*\w+2*\g,0.5*\h) -- (3*\w+3*\g,0.5*\h);
  \node[lbl, anchor=west] at (3*\w+3*\g+0.05,0.5*\h)
    {\tpcifenglish{outcome}{Ergebnis}};
  % Klammer: Stadien 1 und 2 sind zusammen wieder eine Präparation
  \draw[birrengray, line width=0.6pt]
    (0,-0.25) -- (0,-0.45) -- (2*\w+\g,-0.45) -- (2*\w+\g,-0.25);
  \node[lbl, anchor=north] at (\w+0.5*\g,-0.5)
    {\tpcifenglish{together: one more complicated preparation}{zusammen: eine kompliziertere Präparation}};
\end{tikzpicture}
\end{center}

Die ersten beiden Schritte lassen sich zu \emph{einer} komplizierteren Präparation zusammenfassen: Ein gezinkter Würfel ist nichts anderes als ein normaler Würfel, an dem man vorher etwas gemacht hat. Wo man die Grenze zieht, ist Geschmackssache; oft kommt man mit Präparation und Messung aus.
\end{greybox}

Auch die Reihenfolge des Stoffs ist festgelegt: Wir gehen vom Mikroskopischen zum Makroskopischen, von der Quantenmechanik einzelner Systeme über viele Teilchen zur statistischen Mechanik und Thermodynamik. Auf jeder Stufe lernt man, dass jede Theorie eine effektive Theorie der darunterliegenden ist, eine Näherung an die grundlegende Physik darunter.

\subsection{Wahrscheinlichkeit für Physiker}
Die Zahlen, welche die Quantenmechanik als Ergebnis liefert, sind Wahrscheinlichkeiten. Erwarten Sie nichts konzeptionell Neues: Ein intuitives Verständnis setze ich voraus; es geht um die \emph{Notation}, in der wir später die Quantenmechanik formulieren. Wir sprechen dabei die Sprache der Physik: Was die Mathematik eine Wahrscheinlichkeitsverteilung nennt, heißt hier \emph{Präparation} (Zustand), und eine Zufallsvariable heißt \emph{Messung}.

\subsubsection{Ergebnisse und Präparationen}
Um ein Experiment zu modellieren, brauchen wir drei Dinge: die Menge der möglichen Ergebnisse, die Präparation und die Messung. Jedem System wird eine Menge $X$ zugeordnet, die \emph{Ergebnismenge} (Stichprobenraum); man erhält sie, indem man alle möglichen Ergebnisse auflistet. Teilmengen $S \subset X$ heißen \emph{Ereignisse}. Zunächst sei $X$ endlich. Eine \emph{Präparation} wird durch eine Wahrscheinlichkeit $p(x)$ für jedes Ergebnis $x$ charakterisiert, ein \emph{Gewicht} mit
\begin{equation}
	p(x) \geq 0 \quad \text{und} \quad \sum_{x \in X} p(x) = 1.
\end{equation}
Die \emph{Wahrscheinlichkeit} eines Ereignisses ist $\mathbb{P}[S] = \sum_{x \in S} p(x)$.

\begin{beispiel}
	Bei einem Würfel ist $X = \{1, 2, 3, 4, 5, 6\}$, und ,,die Augenzahl ist gerade`` ist das Ereignis $S = \{2, 4, 6\}$. Für einen fairen Würfel ist $p(x) = 1/6$ für alle $x$, also $\mathbb{P}[S] = 1/2$.
\end{beispiel}

\begin{übung}
	Beschreiben Sie die Ergebnismenge für \emph{zwei} Würfel. Wie viele Elemente hat sie?
\end{übung}

Für paarweise disjunkte Ereignisse $S_1, \dots, S_k$ gilt
\begin{equation}
	\mathbb{P}[S_1 \cup \dots \cup S_k] = \mathbb{P}[S_1] + \dots + \mathbb{P}[S_k].
\end{equation}
Die Disjunktheit ist wesentlich: Überlappen sich die Ereignisse, so zählt die rechte Seite Ergebnisse doppelt. Eine solche Mengenfunktion heißt in der Mathematik ein \emph{Maß}, mit $\mathbb{P}[X] = 1$ ein \emph{Wahrscheinlichkeitsmaß}. Mehr Maßtheorie brauchen wir nicht.

\subsubsection{Messungen und Erwartungswert}
Eine \emph{Messung} (Observable, Zufallsvariable) ist eine Funktion $f\colon X \to \mathbb{R}$. Der Wert $f(x)$ ist der Messwert, wenn sicher das Ergebnis $x$ vorliegt. Beachten Sie, wie großzügig das ist: An eine Messung stellen wir keine Bedingung, an die Präparation zwei. Präparation und Messung zusammen ergeben den \emph{Erwartungswert}
\begin{equation}
	\eqbox{\mathbb{E}[f] = \sum_{x \in X} p(x) f(x)}
\end{equation}
Er ist linear, $\mathbb{E}[\alpha f + \gamma g] = \alpha\,\mathbb{E}[f] + \gamma\,\mathbb{E}[g]$, und für eine Konstante ist $\mathbb{E}[c] = c$.

Woher kommt die Formel? Man wiederholt das Experiment $N$-mal mit derselben Präparation und erhält die Ergebnisse $x_1, \dots, x_N$. Kommt $x$ dabei $N(x)$-mal vor, so ist $h_N(x) = N(x)/N$ die \emph{relative Häufigkeit}; sie ist selbst ein Gewicht, und für den Mittelwert der Messreihe gilt exakt
\begin{equation}
	\frac{1}{N} \sum_{i=1}^{N} f(x_i) = \sum_{x \in X} h_N(x)\, f(x).
\end{equation}
Bei vielen unabhängigen Wiederholungen liegt $h_N(x)$ mit hoher Wahrscheinlichkeit nahe bei $p(x)$, und der Mittelwert der Messreihe nahe beim Erwartungswert. Diese Lesart ist eine Deutung: In einer einzelnen Messreihe kann $h_N$ von $p$ abweichen.

\subsubsection{Ja-Nein-Messungen}
Die wichtigsten Messungen nehmen nur die Werte $0$ und $1$ an. Sie sind das Modell für Detektoren, etwa einen Photonendetektor: Klick oder kein Klick. An diesen zweiwertigen Messungen hängt die Quantenmechanik. Eine solche Messung ist durch die Menge $S = \{x \in X \colon f(x) = 1\}$ bestimmt; man schreibt
\begin{equation}
	\chi_S(x) = \begin{cases} 1, & x \in S, \\ 0, & \text{sonst,} \end{cases}
\end{equation}
und nennt $\chi_S$ die \emph{Indikatorfunktion} von $S$. Halten Sie fest: Indikatorfunktionen \emph{sind} Messungen, kein eigenes Objekt daneben. Ihr Erwartungswert ist die Wahrscheinlichkeit ihres Ereignisses:
\begin{equation}\label{eq:indicatorfunction}
	\mathbb{E}[\chi_S] = \sum_{x \in X} p(x)\, \chi_S(x) = \sum_{x \in S} p(x) = \mathbb{P}[S].
\end{equation}
Wahrscheinlichkeiten sind also besondere Erwartungswerte. Umgekehrt legt $\mathbb{P}$ alle Erwartungswerte fest, denn $p(x) = \mathbb{P}[\{x\}]$.

\begin{übung}
	Sei $X = \{1, 2, 3\}$ und $p(x) = 1/3$ für alle $x$. Berechnen Sie $\mathbb{P}[S]$ für $S = \{1, 2\}$ und den Erwartungswert der Messung $f(x) = x^2$.
\end{übung}

\subsubsection{Varianz und Determinismus}
Der Erwartungswert sagt, wo die Messwerte im Mittel liegen; die \emph{Varianz} sagt, wie weit sie streuen:
\begin{equation}
	\operatorname{Var}[f] = \mathbb{E}\bigl[(f - \mathbb{E}[f])^2\bigr] = \mathbb{E}[f^2] - \mathbb{E}[f]^2 .
\end{equation}
Die zweite Form folgt aus der Linearität. Ausgeschrieben ist $\operatorname{Var}[f] = \sum_x p(x)\,(f(x) - \mathbb{E}[f])^2$, eine Summe von Zahlen $\geq 0$. Also ist $\operatorname{Var}[f] \geq 0$, und $\operatorname{Var}[f] = 0$ gilt genau dann, wenn $f(x) = \mathbb{E}[f]$ für alle $x$ mit $p(x) > 0$ ist. In der Quantenmechanik spielt die Varianz eine große Rolle; bei der Unschärferelation begegnet sie uns wieder.

\begin{beispiel}
	Für den fairen Würfel und die Augenzahl $f(x) = x$ ist $\mathbb{E}[f] = 21/6 = 7/2$; der Erwartungswert muss kein möglicher Messwert sein. Mit $\mathbb{E}[f^2] = 91/6$ ist $\operatorname{Var}[f] = 91/6 - 49/4 = 35/12$.
\end{beispiel}

Für eine Ja-Nein-Messung ist $\chi_S^2 = \chi_S$, also
\begin{equation}
	\operatorname{Var}[\chi_S] = \mathbb{P}[S]\,\bigl(1 - \mathbb{P}[S]\bigr).
\end{equation}
Sie verschwindet genau dann, wenn die Antwort feststeht. Ist $p(a) = 1$ für ein Ergebnis $a$, so ist $\mathbb{E}[f] = f(a)$ für jede Messung.

\begin{merksatz}[Determinismus ist ein Sonderfall]
Für ein Gewicht $p$ auf einer endlichen Ergebnismenge sind gleichwertig: $p(a) = 1$ für ein Ergebnis $a$; jedes Ereignis hat die Wahrscheinlichkeit $0$ oder $1$; jede Messung hat die Varianz $0$. Solche Präparationen nennen wir deterministisch: Das Ergebnis jeder Messung steht fest. So passen die definitiven Vorhersagen der klassischen Physik in eine probabilistische Theorie.
\end{merksatz}

\subsubsection{Kontinuierliche Ergebnisse}
Für die Quantenmechanik ist $X = \mathbb{R}$ entscheidend: Ein Teilchen in einer Dimension hat eine Position, und die möglichen Positionen sind die reellen Zahlen. Wir betrachten nur Präparationen mit einer \emph{Wahrscheinlichkeitsdichte}, einer Funktion $p$ auf $\mathbb{R}$ mit
\begin{equation}\label{eq:ctsprob}
	\begin{aligned}
		p(x) &\geq 0, & \int_{-\infty}^{\infty} p(x)\,\dd x &= 1, \\
		\mathbb{P}[[a,b]] &= \int_a^b p(x)\,\dd x, & \mathbb{E}[f] &= \int_{-\infty}^{\infty} p(x) f(x)\,\dd x .
	\end{aligned}
\end{equation}
Dabei ist $a \leq b$, und den Erwartungswert verwenden wir nur, sofern $\int p(x) \abs{f(x)}\,\dd x < \infty$ ist. Gegenüber dem endlichen Fall ist also nur die Summe durch ein Integral ersetzt. In der Literatur steht dafür auch $\int \mathbb{P}(\dd x)\, f(x)$; ich mag diese Schreibweise nicht besonders, Sie sollen sie nur erkennen. Wichtig: Der Wert $p(x)$ einer Dichte ist keine Wahrscheinlichkeit. Wahrscheinlichkeiten sind Flächen unter der Kurve; ein einzelner Punkt hat die Wahrscheinlichkeit $0$, und $p(x)$ darf größer als $1$ sein.

\begin{beispiel}
	Die Dichte
	\begin{equation}
		p(x) = \frac{1}{\sqrt{2\pi}}\, e^{-\frac12 x^2}
	\end{equation}
	heißt \emph{Normal-} oder \emph{Gauß-Verteilung}; sie ist normiert. Ist das Teilchen links oder rechts vom Ursprung? Weil $p(-x) = p(x)$ ist, sind die Flächen links und rechts vom Nullpunkt gleich groß, zusammen $1$; also ist $\int_0^\infty p(x)\,\dd x = 1/2$. Vorsicht: Das Argument funktioniert \emph{nur} für symmetrische Dichten.
\end{beispiel}

Die Erwartungswerte der Koordinatenfunktion und ihrer Potenzen heißen \emph{Momente}, sofern die Integrale $\int \abs{x}^n p(x)\,\dd x$ endlich sind:
\begin{equation}
	M_n = \int_{-\infty}^{\infty} x^n\, p(x)\,\dd x .
\end{equation}
Nur $M_0 = 1$, die Normierung, existiert immer; $M_1$ heißt, sofern es existiert, \emph{Mittelwert}. Existieren $M_1$ und $M_2$, so ist die Varianz der Koordinate wie im endlichen Fall
\begin{equation}
	\operatorname{Var}[\mathbb{P}] = \int_{-\infty}^{\infty} (x - M_1)^2\, p(x)\,\dd x = M_2 - M_1^2 .
\end{equation}
Mittelwert und Varianz tragen dann dieselbe Information wie $M_1$ und $M_2$. Für die hier betrachteten Modelle steht damit die Sprache bereit: Präparationen sind Gewichte oder Dichten, Messungen sind Funktionen, und beide zusammen ergeben bei endlichen Integralen Erwartungswerte.
